% Source: 11 - Tue Oct 20 2026.pdf, page 11. Content remains in source order.
\sourcecontinuation{11}{11}
considering the sequence of active calls of procedure PREORDER($T,r$).
\end{notesdefinition}
\textbf{EXAMPLE:}
\sourcefigure[0.48]{assets/11/page-011-preorder-walk.png}
\begin{tabular}{@{}ll@{}}
PREORDER($T,a$) calls & PREORDER($T,b$)\\
PREORDER($T,b$) calls & PREORDER($T,c$)\\
PREORDER($T,c$) returns to & PREORDER($T,b$)\\
PREORDER($T,b$) calls & PREORDER($T,d$)\\
PREORDER($T,d$) returns to & PREORDER($T,b$)\\
PREORDER($T,b$) returns to & PREORDER($T,a$)\\
PREORDER($T,a$) calls & PREORDER($T,e$)\\
PREORDER($T,e$) returns to & PREORDER($T,a$)
\end{tabular}
\[\mathrm{FPW}=\langle a,b,c,b,d,b,a,e,a\rangle.\]
\begin{notesresult}
\textbf{PROPERTIES OF A FPW}
\begin{enumerate}
\item The first occurrence of each node gives the preorder visit of $T$.
\item Each edge of $T$ is used 2 times by FPW:
\end{enumerate}
\end{notesresult}
\sourcefigure[0.72]{assets/11/page-011-edge-traversals.png}
Therefore:
\[c(\mathrm{FPW})=2c(T^*).\]
